\chapter{Divisores: certificar uma multiplicação de pontos}
\label{chapter:divisors}
\label{sec:fcmp-divisors}

\section{A pergunta que uma função pode responder}

A afirmação \(Q=nP\) tem um sentido conhecido: somar \(P\) consigo
mesmo \(n\) vezes produz \(Q\). Podemos confirmá-la por duplicações e
adições. Que objecto poderia acompanhar \(Q\) e certificar o resultado?

Uma função na curva pode desempenhar esse papel. A ponte está nos seus
zeros e pólos: os lugares onde ocorrem, contados com multiplicidade,
obedecem a uma relação na lei do grupo. Construiremos essa ligação desde
uma recta até uma função que testemunha a multiplicação.

Para estudar essa ponte, usaremos uma curva não singular na forma curta de
Weierstrass, sobre um corpo \(\mathbb F\) de característica maior que três:
\begin{equation}
 \mathbb E:\quad y^2=x^3+a_{\mathbb E}x+b_{\mathbb E},
 \qquad 4a_{\mathbb E}^{\,3}+27b_{\mathbb E}^{\,2}\ne0.
 \label{eq:fcmp-short-weierstrass}
\end{equation}
Chamaremos \(\mathcal O_\infty\) à identidade. Nesta apresentação,
\(-(x,y)=(x,-y)\). O inverso em coordenadas de Edwards, estabelecido no
capítulo anterior, tem outra fórmula: a representação dos pontos importa.

O nosso exemplo será deliberadamente pequeno:
\[
 \mathbb E/\mathbb F_{17}:\quad y^2=x^3+2x+2,
 \qquad P=(5,1),\qquad Q=(10,6).
\]
As contas deste exemplo serão módulo \(17\). Construiremos a função que
certifica \(Q=3P\) e examinaremos o seu divisor. Um corpo pequeno permite
conferir as contas; a solidez criptográfica exigirá corpos grandes.

\newpage
\section{Uma recta, três encontros e uma soma}

Na forma de Weierstrass, a soma tem uma leitura geométrica. Uma recta passa
por \(A\) e \(B\), encontra a curva num terceiro ponto \(R\), e a soma
\(A+B\) é o inverso de \(R\). Se os dois pontos coincidem, usa-se a tangente.
Se são inversos, a recta é vertical e a soma é \(\mathcal O_\infty\).

Para uma recta não vertical \(y=mx+c\), a substituição na curva dá
\[
 x^3-m^2x^2+(a_{\mathbb E}-2mc)x+b_{\mathbb E}-c^2=0.
\]
O polinómio é cúbico. As suas raízes são as coordenadas dos três encontros,
contados com multiplicidade, numa extensão do corpo se necessário. Pela soma
das raízes, \(x_A+x_B+x_R=m^2\); assim,
\[
 x_{A+B}=m^2-x_A-x_B,
 \qquad y_{A+B}=m(x_A-x_{A+B})-y_A.
\]
São operações no corpo, mesmo quando a figura que as sugere é desenhada sobre
os reais. Num corpo finito, a curva é um conjunto de pontos discretos; a
equação de intersecção conserva o raciocínio algébrico.

No exemplo, a tangente em \(P=(5,1)\) tem declive
\[
 m=\frac{3\cdot5^2+2}{2\cdot1}
   =\frac9{2}=13\pmod {17}.
\]
Logo, \(2P=(6,3)\). A recta por \(P\) e \(2P\) tem declive
\((3-1)/(6-5)=2\), e fornece \(3P=(10,6)=Q\).

\begin{center}
\begin{tikzpicture}[>=stealth, node distance=1.1cm]
 \node[draw,rounded corners,align=center] (a) {\(A,B\)\\pontos dados};
 \node[draw,rounded corners,align=center,right=of a] (b)
 {recta\\terceiro encontro \(R\)};
 \node[draw,rounded corners,align=center,right=of b] (c)
 {reflexão\\\(A+B=-R\)};
 \draw[->] (a)--(b); \draw[->] (b)--(c);
\end{tikzpicture}
\end{center}

A relação que interessa guardar é \(A+B+R=\mathcal O_\infty\).
Cada recta oferece, portanto, uma pequena soma nula. O certificado final
reunirá estas relações e eliminará os pontos intermédios.

\newpage
\section{A função existe sobre a curva}

A expressão \(\lambda(x,y)=y-mx-c\) define uma função. Avaliá-la num ponto
da curva dá \(\lambda(A)=y_A-mx_A-c\). O resultado é zero exactamente
quando \(A\) pertence à recta. A recta desenhada no plano e a função
restringida à curva têm papéis distintos: a primeira indica os encontros;
a segunda permite multiplicar, dividir e comparar valores.

Na curva, as coordenadas não são independentes. A função \(y^2\) é a mesma
que \(x^3+a_{\mathbb E}x+b_{\mathbb E}\). Por isso, o anel das funções
polinomiais na parte afim é
\[
 \mathbb F[\mathbb E]=
 \mathbb F[x,y]/(y^2-x^3-a_{\mathbb E}x-b_{\mathbb E}).
\]
O quociente diz que identificamos expressões cuja diferença é um múltiplo
da equação da curva. Reduzindo as potências de \(y\), cada elemento admite
um representante único
\begin{equation}
 f(x,y)=f_0(x)+yf_1(x),\qquad f_0,f_1\in\mathbb F[x].
 \label{eq:fcmp-reduced-function-representative}
\end{equation}
Por exemplo, \(y^3= y(x^3+a_{\mathbb E}x+b_{\mathbb E})\).
Esta redução tornará possível escrever uma função como uma lista finita de
coeficientes, adequada a uma testemunha algébrica.

O corpo de funções \(\mathbb F(\mathbb E)\) é o corpo de fracções deste
anel: os seus elementos são razões de funções polinomiais, com denominador
não identicamente nulo. Uma razão pode não estar definida pela sua expressão
num ponto particular; é preciso determinar se há cancelamento ou um pólo.

Há uma precaução própria dos corpos finitos. Consideramos funções como
objectos algébricos na curva, e não apenas tabelas dos valores nos seus poucos
pontos racionais. Um polinómio pode anular-se em todos esses pontos sem ser
a função identicamente nula. Os zeros geométricos incluem também os pontos
definidos sobre extensões do corpo.

\textbf{Verificação de compreensão.} Que função representa
\(y^2-x^3-2x-2\) no nosso exemplo? A função nula. Já \(y-2x+9\) é uma
função não nula, embora se anule em três encontros com a curva. Um certificado
precisará sempre de uma função não nula.

\newpage
\section{Um zero pode contar mais de uma vez}

Saber que \(f(A)=0\) ainda não diz quantas vezes esse zero conta. Num
polinómio de uma variável, \((t-t_0)^2\) tem um zero duplo em \(t_0\).
Na curva, procede-se localmente: escolhe-se uma coordenada que se anule
simplesmente no ponto e mede-se a primeira potência não nula da função.

Essa coordenada chama-se \emph{parâmetro local}. Se \(y_A\ne0\), podemos
usar \(t=x-x_A\); a equação da curva determina uma expansão local de
\(y\) em potências de \(t\). Uma função não nula escreve-se
\[
 f=t^m u,\qquad u(A)\ne0.
\]
O inteiro \(m\) é a ordem de \(f\) em \(A\), escrita
\(\operatorname{ord}_A(f)\). Uma ordem positiva indica um zero; uma ordem
negativa, um pólo; uma ordem nula, um valor local finito e não nulo.

A tangente elimina o termo linear da expansão. Num ponto comum, a sua
função tem, por isso, um zero de ordem dois. Num ponto de inflexão, a ordem
pode ser três. O número de pontos distintos e o número de intersecções
contadas com multiplicidade deixam então de coincidir.

No nosso exemplo, a tangente é
\[
 \lambda_{P,P}=y-13x+13.
\]
Substituir \(y=13x-13\) na equação da curva produz o polinómio mónico
\[
 (x-5)^2(x-6).
\]
Há duas ocorrências em \(P=(5,1)\) e uma em \(-2P=(6,14)\). A tangente
não precisa de encontrar três pontos distintos para fornecer três zeros
contados correctamente.

Se \(y_A=0\), a expressão \(x-x_A\) já não é, em geral, um parâmetro
local: neste caso, \(y\) serve como parâmetro e \(x-x_A\) tem ordem dois.
É por isso que a ordem pertence à função na curva, e não apenas à aparência
da sua expressão no plano.

\textbf{Verificação de compreensão.} Se \(f\) tem um zero duplo e \(g\)
um zero simples no mesmo ponto, qual é a ordem de \(f/g\)? É um:
localmente, \(t^2u/(tv)=t(u/v)\). O cancelamento tem de respeitar as ordens.

\newpage
\section{O infinito entra na contabilidade}

Uma função polinomial não tem pólos na parte afim da curva. Isso não impede
que os tenha no ponto \(\mathcal O_\infty\). Para ver esse ponto,
homogeneizamos a equação:
\[
 Y^2Z=X^3+a_{\mathbb E}XZ^2+b_{\mathbb E}Z^3.
\]
As coordenadas afins são \(x=X/Z\), \(y=Y/Z\). Quando \(Z=0\), a equação
obriga a \(X=0\), ficando o único ponto \([0:1:0]\). É a identidade da
curva nesta representação.

Na carta \(Y\ne0\), ponhamos \(u=X/Y\) e \(v=Z/Y\). A equação torna-se
\[
 v=u^3+a_{\mathbb E}uv^2+b_{\mathbb E}v^3.
\]
Perto da identidade, \(u\) é um parâmetro local e \(v\) começa em ordem
três. Como \(x=u/v\) e \(y=1/v\), obtemos
\[
 \operatorname{ord}_{\mathcal O_\infty}(x)=-2,
 \qquad \operatorname{ord}_{\mathcal O_\infty}(y)=-3.
\]
São ordens algébricas; não estamos a fazer tender números de um corpo finito
para uma grandeza infinita.

A função vertical \(x-x_A\) tem, portanto, um pólo duplo na identidade.
A função de uma recta não vertical, \(y-mx-c\), tem um pólo triplo:
o termo \(y\), de ordem \(-3\), não pode ser cancelado por termos de
ordens \(-2\) e zero. Os dois zeros afins da vertical e os três zeros
afins da recta equilibram esses pólos.

Também se explica a escolha dos pesos \(2\) para \(x\) e \(3\) para
\(y\). No representante reduzido \(f_0+yf_1\), os monómios \(x^i\)
têm ordem de pólo \(2i\), e os monómios \(yx^j\), ordem \(3+2j\).
Um peso é par e o outro ímpar; os termos de maior peso não se cancelam entre
as duas parcelas. Este peso mede a ordem do pólo no infinito.

O ponto no infinito é, pois, indispensável: deixa visível o saldo que uma
função polinomial transporta para fora da parte afim.

\newpage
\section{O divisor: um registo de lugares e ordens}

Podemos agora dar nome ao registo que temos vindo a preparar. Um
\emph{divisor} é uma soma formal finita
\begin{equation}
 D=\sum_{A\in S}m_A[A],\qquad m_A\in\mathbb Z.
 \label{eq:fcmp-formal-divisor}
\end{equation}
O símbolo \([A]\) guarda o lugar \(A\); o coeficiente guarda uma
multiplicidade inteira. Neste capítulo, os divisores serão escritos sobre
os pontos geométricos, isto é, sobre o fecho algébrico do corpo. Quando todos
os pontos usados são racionais, como no exemplo, as contas já se fazem
sobre \(\mathbb F\).

Somar divisores é somar coeficientes do mesmo símbolo. Assim,
\[
 2[A]+[B]-[A]=[A]+[B].
\]
Esta operação não soma pontos na curva. Mesmo que \(C=A+B\), o divisor
\([A]+[B]\) conserva dois lugares, enquanto \([C]\) conserva um só.
O grupo dos divisores é o grupo abeliano livre gerado pelos símbolos dos
pontos; nenhuma relação de adição da curva é imposta nesta definição.

O \emph{suporte} de \(D\) é o conjunto dos lugares com coeficiente não
nulo. O seu \emph{grau} é
\[
 \deg D=\sum_A m_A.
\]
Por exemplo, \(2[P]+[-2P]-3[\mathcal O_\infty]\) tem grau zero e, no
nosso exemplo, suporte com três pontos distintos. Coeficientes negativos
pertencem à definição de divisor; só quando o divisor provém de uma função
significam efectivamente ordens de pólos.

Para uma função racional não nula, o divisor é
\begin{equation}
 \operatorname{div}(f)=\sum_A\operatorname{ord}_A(f)[A].
 \label{eq:fcmp-rational-function-divisor}
\end{equation}
Chamamos \emph{principal} a um divisor que tem esta forma. Uma constante
não nula tem divisor zero. A função identicamente nula não tem divisor:
não admite uma ordem finita nos seus pontos.

\textbf{Verificação de compreensão.} Quanto vale
\([A]+[-A]\)? Como divisor, continua a ser essa soma formal, salvo se os
lugares coincidirem. A soma dos pontos \(A+(-A)\), por sua vez, é a
identidade. Os parênteses rectos preservam a distinção.

\newpage
\section{Multiplicar funções soma os seus divisores}

Se localmente \(f=t^m u\) e \(g=t^n v\), com \(u,v\) unidades, então
\(fg=t^{m+n}uv\). A ordem de um produto é a soma das ordens. Seguem-se
as identidades
\[
 \operatorname{div}(fg)=\operatorname{div}(f)+\operatorname{div}(g),
 \qquad
 \operatorname{div}(f/g)=\operatorname{div}(f)-\operatorname{div}(g).
\]
Esta é a primeira ferramenta de construção: relações pequenas podem ser
reunidas por multiplicação; os lugares intermédios podem desaparecer por
divisão.

Consideremos \(C=A+B\), em posição genérica. A recta por \(A,B\) tem
terceiro encontro em \(-C\). A vertical por \(C\) passa também por \(-C\).
Os respectivos divisores são
\begin{align*}
 \operatorname{div}(\lambda_{A,B})
   &=[A]+[B]+[-C]-3[\mathcal O_\infty],\\
 \operatorname{div}(v_C)
   &=[C]+[-C]-2[\mathcal O_\infty],
 \qquad v_C=x-x_C.
\end{align*}
Ao dividir, o encontro \(-C\) cancela-se e resta
\begin{equation}
 h_{A,B}=\frac{\lambda_{A,B}}{v_C},\qquad
 \operatorname{div}(h_{A,B})=[A]+[B]-[C]-[\mathcal O_\infty].
 \label{eq:fcmp-small-principal-divisor}
\end{equation}
No caso genérico, os zeros são \(A,B\) e os pólos são
\(C,\mathcal O_\infty\). Se houver coincidências, somam-se primeiro os
coeficientes: a lista de zeros e pólos efectiva pode ser menor.

Para duplicar \(P\), a mesma construção dá
\[
 \operatorname{div}(h_{P,P})=2[P]-[2P]-[\mathcal O_\infty].
\]
A tangência já fornece o coeficiente dois; não se acrescenta uma segunda
condição de avaliação num ponto diferente.

O divisor da razão tem grau zero. Esta igualdade reflecte um facto geral:
numa curva projectiva não singular, toda a função racional não nula tem
tantos zeros como pólos, contando multiplicidades. A condição é necessária
para um divisor ser principal. Precisamos, contudo, de uma segunda condição
que guarde a posição dos lugares, pois uma mera contagem perde a relação
entre os pontos.

\newpage
\section{Duas somas que convém nunca confundir}

A cada divisor podemos associar dois resultados. O primeiro é o inteiro
\(\deg D\). O segundo é um ponto, obtido aplicando a lei do grupo aos
lugares com as suas multiplicidades:
\[
 \sigma(D)=\sum_A m_A A.
\]
Coeficientes negativos mandam somar inversos; coeficientes positivos mandam
somar repetidamente. O símbolo \(D\) é o registo formal; \(\sigma(D)\)
é o resultado da soma dos pontos registados.

Para o divisor da recta por \(A,B\), temos
\[
 \deg D=3-3=0,
 \qquad \sigma(D)=A+B+(-A-B)=\mathcal O_\infty.
\]
Para a vertical, a soma é \(A+(-A)=\mathcal O_\infty\). Produtos e
quocientes destas funções conservam as duas condições. Tanto o grau como
\(\sigma\) são aditivos.

O exemplo \(D=[P]-[Q]\) mostra a diferença entre as condições. O grau
é zero para quaisquer \(P,Q\), mas \(\sigma(D)=P-Q\) só é a identidade
quando \(P=Q\). Se \(P\ne Q\), não existe uma função com um único zero
simples em \(P\) e um único pólo simples em \(Q\). A curva elíptica impõe
mais do que o equilíbrio do número de zeros e pólos.

O teorema que usaremos é precisamente
\begin{equation}
 D\text{ principal}\quad\Longleftrightarrow\quad
 \deg D=0\ \text{ e }\ \sigma(D)=\mathcal O_\infty.
 \label{eq:fcmp-principal-divisor-criterion}
\end{equation}
A formulação vale para divisores geométricos na curva elíptica. Para os
divisores de pontos racionais considerados aqui, a função testemunha pode
ser tomada sobre o corpo base. A construção por rectas que faremos fornece
essa função directamente.

Este critério é a ponte procurada. Um registo formal cuja soma no grupo é
zero pode ser o divisor de uma função; reciprocamente, o divisor de uma
função tem soma nula. Para compreender a afirmação, falta explicar por que
a geometria das rectas representa todas as classes de grau zero. Esse será
o próximo passo.

\newpage
\section{Por que a ponte respeita a lei do grupo}

Dois divisores dizem-se \emph{linearmente equivalentes} quando a sua
diferença é principal. Escrevemos \(D\sim D'\). A função
\(h_{A,B}\) construída anteriormente demonstra
\[
 [A]+[B]\sim[A+B]+[\mathcal O_\infty].
\]
Os divisores dos dois lados continuam diferentes. O sinal \(\sim\) diz
que existe uma função cuja contabilidade transforma um no outro.

Subtraindo duas cópias da identidade, obtemos
\[
 ([A]-[\mathcal O_\infty])+([B]-[\mathcal O_\infty])
 \sim[A+B]-[\mathcal O_\infty].
\]
Esta relação deixa reduzir uma soma de diferenças a uma única diferença.
As verticais permitem tratar os inversos, e a identidade não acrescenta
nada. Para qualquer divisor de grau zero, a redução conduz a
\[
 D\sim[\sigma(D)]-[\mathcal O_\infty].
\]
Em particular, se \(\sigma(D)=\mathcal O_\infty\), o representante final
é zero. As funções usadas nas reduções, multiplicadas ou divididas conforme
o caso, fornecem uma função de divisor \(D\).

Falta justificar a unicidade desse representante. Se
\([S]-[\mathcal O_\infty]\) fosse principal com \(S\ne\mathcal O_\infty\),
a função correspondente teria um único pólo simples, na identidade, e
nenhum pólo afim. Seria, pois, uma função polinomial na parte afim. Mas o
representante \(f_0+yf_1\) mostra que uma função polinomial não constante
tem ordem de pólo pelo menos dois: não há monómio de peso um. Tal função
não existe. Isto completa a justificação do critério.

O conjunto de classes de divisores de grau zero chama-se
\(\operatorname{Pic}^0(\mathbb E)\). A correspondência
\[
 S\longmapsto\text{classe de }[S]-[\mathcal O_\infty]
\]
é um isomorfismo com o grupo da curva. O nome é útil, mas o mecanismo já
está à vista: uma recta transporta a adição dos pontos para a equivalência
dos divisores. É esse mecanismo que o certificado aproveitará.

\newpage
\section{Construir o certificado de uma multiplicação}

Voltemos à relação \(Q=nP\), com \(n\ge1\). O divisor
\[
 D_n=n[P]-[Q]-(n-1)[\mathcal O_\infty]
\]
tem grau zero. A sua soma no grupo é \(nP-Q\). Pelo critério, existe
uma função com este divisor exactamente quando \(Q=nP\).

A existência pode ser acompanhada passo a passo. Ponhamos \(F_1=1\).
Supondo \(Q_k=kP\) e
\[
 \operatorname{div}(F_k)=k[P]-[Q_k]-(k-1)[\mathcal O_\infty],
\]
a função da próxima adição satisfaz
\[
 \operatorname{div}(h_{Q_k,P})
 =[Q_k]+[P]-[Q_{k+1}]-[\mathcal O_\infty].
\]
Logo, \(F_{k+1}=F_kh_{Q_k,P}\) tem o divisor desejado. O termo
\(-[Q_k]\) cancela-se com \([Q_k]\). A função final guarda a relação
entre os extremos sem conservar os pontos intermédios no seu divisor.

Também podemos seguir uma cadeia de duplicações e adições. Se já temos
\(F_u,F_v\), então
\[
 F_{u+v}=F_uF_vh_{uP,vP}
\]
tem divisor \((u+v)[P]-[(u+v)P]-(u+v-1)[\mathcal O_\infty]\).
Para duplicar, usa-se \(F_u^2h_{uP,uP}\). A construção deixa assim
acompanhar o algoritmo habitual de multiplicação escalar.

Para \(n=3\), a conta formal inteira cabe em duas linhas:
\begin{align*}
 \operatorname{div}(h_{P,P})&=2[P]-[2P]-[\mathcal O_\infty],\\
 \operatorname{div}(h_{2P,P})&=[2P]+[P]-[3P]-[\mathcal O_\infty].
\end{align*}
A soma dá \(3[P]-[3P]-2[\mathcal O_\infty]\).

Há aqui uma diferença de custos a conservar. O provador ainda calcula os
pontos e constrói a função. O certificado abre uma possibilidade de
verificação diferente; não elimina o trabalho necessário para obter a
testemunha. Além disso, a ordem \(n\) do zero em \(P\) seria enorme para
um escalar criptográfico. Teremos de escolher uma representação mais económica.

\newpage
\section{O certificado completo em módulo dezassete}

No exemplo, já encontrámos \(2P=(6,3)\) e \(Q=3P=(10,6)\). As duas
rectas e as duas verticais são
\begin{align*}
 \lambda_{P,P}&=y-13x+13,&v_{2P}&=x-6,\\
 \lambda_{2P,P}&=y-2x+9,&v_Q&=x-10.
\end{align*}
Assim,
\[
 F_3=\frac{(y-13x+13)(y-2x+9)}{(x-6)(x-10)}
\]
tem divisor \(3[P]-[Q]-2[\mathcal O_\infty]\). O cancelamento em
\(2P\) exprime-se nas ordens locais, mesmo que a fracção escrita
pareça indeterminada nesse ponto.

Para preparar uma testemunha polinomial, multiplicamos por \(v_Q\). A
função \(f=F_3v_Q\) fica com divisor
\[
 \operatorname{div}(f)=3[P]+[-Q]-4[\mathcal O_\infty].
\]
Reduzir o produto das rectas pela equação da curva dá
\[
 (y-13x+13)(y-2x+9)
 =(x-6)\bigl(x^2+15x+2y\bigr).
\]
Podemos, portanto, tomar
\[
 f=x^2+15x+2y.
\]
O maior peso é quatro, vindo de \(x^2\). A função tem um pólo de ordem
quatro na identidade. Em \(P\), a construção pelas rectas fornece um zero
triplo. Em \(-Q=(10,11)\), fornece um zero simples. São quatro zeros:
não há espaço para outros.

Uma substituição ordinária confirma
\(f(5,1)=0\) e \(f(10,11)=0\), mas não mede, por si só, o zero triplo.
Essa multiplicidade vem da tangente e do cancelamento local. Pelo contrário,
\(f(10,6)=7\ne0\): o zero adicional está em \(-Q\), não em \(Q\).

O certificado tem agora uma forma concreta: três coeficientes não nulos
e uma ordem de pólo limitada. O divisor tem soma \(3P-Q\), logo a
existência desta função com estas ordens certifica a multiplicação. A
próxima simplificação reduzirá também o número de ocorrências de \(P\).

\newpage
\section{Os dígitos substituem uma multiplicidade enorme}

Para um escalar \(s_A\), escrevamos uma decomposição binária e preparemos
uma tabela pública de pontos:
\begin{equation}
 s_A=\sum_{i=0}^{\ell_A-1}e_i2^i,
 \qquad e_i\in\{0,1\},\qquad A_i=2^iA.
 \label{eq:fcmp-scalar-decomposition}
\end{equation}
A relação a certificar é
\begin{equation}
 W=s_AA\quad\Longleftrightarrow\quad
 \sum_i e_iA_i-W=\mathcal O_\infty.
 \label{eq:fcmp-hidden-scalar-multiplication}
\end{equation}
Em lugar de pôr \(s_A\) ocorrências em \(A\), usamos uma ocorrência em
cada ponto seleccionado pela decomposição. O divisor correspondente é
\begin{equation}
 D_{s_A}=\sum_i e_i[A_i]+[-W]
       -\left(1+\sum_i e_i\right)[\mathcal O_\infty].
 \label{eq:fcmp-dlog-divisor}
\end{equation}
O grau é zero, e a soma no grupo é \(s_AA-W\). O número de zeros afins
é, no máximo, \(\ell_A+1\), em vez de crescer com o valor inteiro do
escalar. Como só há pólos no infinito, a função testemunha pode ser
representada por um polinómio na curva.

No caso \(3=1+2\), a tabela selecciona \(P\) e \(2P\). O divisor passa
a ser \([P]+[2P]+[-Q]-3[\mathcal O_\infty]\), e a própria recta
\(y-2x+9\) é uma testemunha. A função de peso quatro da página anterior
era correcta; a decomposição permitiu uma de peso três.

Os bits precisam de ser realmente bits, por exemplo pela condição
\(e_i(e_i-1)=0\). A recomposição tem de usar os mesmos dígitos que o
divisor. Se a representação deve ser canónica, impõe-se ainda o intervalo
do escalar. As contas dos coeficientes usam o corpo base; as contas dos
escalares seguem a ordem do subgrupo. Estes módulos não são intercambiáveis.

Se \(W=\mathcal O_\infty\), o termo \([-W]\) junta-se ao termo do
infinito; não se inventam coordenadas afins para a identidade. Se
\(s_A=0\) e \(W=\mathcal O_\infty\), o divisor é zero e uma constante
não nula serve de testemunha. Casos de identidade exigem esse tratamento
explícito nas representações e nos dispositivos de verificação.

\newpage
\section{O tamanho da função também faz parte da prova}

Uma função pode anular-se numa lista qualquer de pontos se lhe concedermos
grau suficiente. Por exemplo, o produto das verticais desses pontos
anula-se também nos respectivos inversos. A mera existência de zeros nos
lugares declarados não demonstra que a soma dos pontos seja a identidade.
É preciso excluir zeros adicionais.

Se o divisor prescrito tem \(d\) zeros afins, contando multiplicidades,
uma função polinomial testemunha deve ter pólo de ordem \(d\) no infinito.
O representante reduzido impõe limites transparentes:
\[
 2\deg f_0\le d,
 \qquad 3+2\deg f_1\le d.
\]
A segunda condição desaparece quando \(f_1=0\). Estes limites dão uma
ordem de pólo de, no máximo, \(d\). Se já houver \(d\) zeros prescritos
com as multiplicidades certas, o equilíbrio do divisor obriga a que todos
os zeros tenham sido contabilizados.

Esta observação fornece uma verificação exacta, embora potencialmente
dispendiosa: confirmar que \(f\ne0\), confirmar o limite de peso e
confirmar cada zero com a respectiva ordem. Para lugares repetidos, avaliar
\(f\) uma vez não basta. Verifica-se a expansão local ou condições de
interpolação que expressem a multiplicidade. Em característica positiva,
derivadas ordinárias de ordem elevada podem perder informação; coeficientes
da expansão local evitam essa ambiguidade.

As funções \(f_0+yf_1\) de peso limitado formam um espaço vectorial de
dimensão aproximadamente \(d\), e exactamente \(d\) para \(d\ge1\).
Para \(d=3\), a base é \(1,x,y\); para \(d=4\), é
\(1,x,x^2,y\). As condições de zeros constituem equações lineares nos
coeficientes, quando os lugares e as condições locais são conhecidos.

No exemplo binário de \(3P\), o peso três força uma recta. Uma recta por
\(P\) e \(2P\) só admite como terceiro encontro \(-Q\). O limite da
função transforma a interpolação numa afirmação sobre a soma. Num protocolo,
esta ligação entre limites, multiplicidades e pontos tem de ser conservada.

\newpage
\section{Avaliar uma função num divisor}

Podemos substituir muitas avaliações individuais por uma identidade entre
produtos. Para um divisor \(D=\sum_A m_A[A]\), define-se
\[
 f(D)=\prod_A f(A)^{m_A},
\]
desde que \(f\) não tenha zeros nem pólos no suporte de \(D\). Os
expoentes negativos significam inversas multiplicativas. A expressão
depende das multiplicidades, tal como o divisor.

A reciprocidade de Weil afirma, para funções racionais não nulas com
divisores de suportes disjuntos,
\[
 f(\operatorname{div}g)=g(\operatorname{div}f).
\]
Do lado esquerdo, avaliamos \(f\) nos zeros e pólos de \(g\). Do lado
direito, avaliamos \(g\) nos zeros e pólos de \(f\). A identidade permite
escolher \(g\) simples e deslocar para essa função as numerosas avaliações
da lista que queremos certificar.

Para visualizar a troca, suponhamos que os zeros de \(f\) são
\(P_1,\ldots,P_d\) e que \(g\) tem apenas três zeros. Avaliar uma
função complicada em três lugares pode ser mais económico do que avaliá-la
nos \(d\) lugares. Do outro lado ficam \(d\) avaliações de uma função
simples. A reciprocidade não elimina essas parcelas; muda a sua natureza.

Há, contudo, um obstáculo que a fórmula curta não resolve. A nossa função
\(f\) é polinomial na curva, e queremos que \(g\) seja a função de uma
recta. Ambas têm um pólo em \(\mathcal O_\infty\). Os suportes não são
disjuntos. Substituir ingenuamente os valores no infinito produziria
expressões sem sentido.

Precisamos de conservar a contribuição local nesse ponto antes de cancelar
as ordens. A regularização da próxima página fará precisamente isso. Esta
precaução é necessária para obter uma identidade exacta, com o sinal e os
coeficientes correctos. Uma igualdade vaga de produtos esconderia a parte
da contabilidade que já aprendemos a não omitir.

\newpage
\section{A contribuição local e o resultante elíptico}

Se num lugar \(A\) escrevemos \(f=t^{m_f}u_f\) e \(g=t^{m_g}u_g\),
as unidades \(u_f,u_g\) têm valores finitos e não nulos. O símbolo local
guarda o valor depois do cancelamento das potências do parâmetro:
\begin{equation}
 \{f,g\}_A=(-1)^{m_fm_g}
       \frac{u_f(A)^{m_g}}{u_g(A)^{m_f}},
 \qquad \prod_A\{f,g\}_A=1.
 \label{eq:fcmp-local-symbol}
\end{equation}
A segunda igualdade é a forma geral da reciprocidade. Quando os suportes
são disjuntos, recupera a identidade da página anterior. Quando coincidem,
permite tratar a contribuição comum sem avaliar um zero ou um pólo como
se fosse um valor ordinário.

Para polinómios na curva, ponhamos
\(\deg_\infty(f)=-\operatorname{ord}_{\mathcal O_\infty}(f)\).
O coeficiente do monómio de maior peso no representante reduzido chama-se
\(\operatorname{lc}(f)\). Escolhendo \(t=x/y\) no infinito, os termos
líderes de \(x\) e \(y\) são \(t^{-2}\) e \(t^{-3}\), com coeficiente
um. Esta escolha torna explícita a correspondência com o peso elíptico.

Se os suportes finitos de \(f,g\) forem disjuntos, o resultante elíptico
admite as duas expressões
\begin{align}
 \operatorname{Res}_{\mathbb E}(f,g)
 &=\operatorname{lc}(g)^{\deg_\infty(f)}
   \prod_{\substack{g(A)=0\\A\ne\mathcal O_\infty}}
      f(A)^{\operatorname{ord}_A(g)}\nonumber\\
 &=(-1)^{\deg_\infty(f)\deg_\infty(g)}
   \operatorname{lc}(f)^{\deg_\infty(g)}
   \prod_{\substack{f(A)=0\\A\ne\mathcal O_\infty}}
      g(A)^{\operatorname{ord}_A(f)}.
 \label{eq:fcmp-weil-reciprocity}
\end{align}
O sinal e os coeficientes líderes são a contribuição do infinito. Esta
forma permite a troca de avaliações que procurávamos, respeitando o pólo
comum. O resultante pode ainda ser definido quando há zeros comuns, mas
as razões que usaremos exigem desafios fora desses casos degenerados.

É nesta identidade regularizada que assenta o método de certificação por
avaliações desafiadas \cite{eagen-divisors}. O teorema dos divisores certifica a soma; o resultante
oferece uma maneira de testar a função testemunha contra a lista de pontos.

\newpage
\section{Uma recta desafiada interroga a testemunha}

Depois de fixar a função testemunha, escolhemos a recta
\begin{equation}
 \mathcal L(x,y)=y-c_{\rm declive}x-c_{\rm corte},
 \qquad Z(x,y)=y-c_{\rm declive}x.
 \label{eq:fcmp-random-line}
\end{equation}
Chamemos \(B_0,B_1,B_2\) aos seus encontros com a curva. Para usar
avaliações no corpo base, o procedimento de desafio tem de fornecer esses
pontos no corpo. Uma construção possível escolhe dois pontos distintos,
não inversos, e toma o terceiro encontro da recta por eles. Escolher dois
coeficientes uniformes não garante que os três encontros sejam racionais.

Suponhamos encontros simples e nenhum zero comum finito com \(f\).
Se \(\operatorname{div}(f)=\sum_j m_j[P_j]-d[\mathcal O_\infty]\),
com \(m_j>0\) e \(\sum_jm_j=d\), o resultante dá
\[
 \prod_{r=0}^2 f(B_r)
 =(-1)^{3d}\operatorname{lc}(f)^3
    \prod_j\mathcal L(P_j)^{m_j}.
\]
Como \(\mathcal L(P_j)=Z(P_j)-c_{\rm corte}\), os sinais cancelam-se:
\[
 \prod_{r=0}^2 f(B_r)
 =\operatorname{lc}(f)^3
   \prod_j(c_{\rm corte}-Z(P_j))^{m_j}.
\]
A função trabalhosa aparece em três avaliações. A lista dos zeros aparece
através de avaliações da função linear \(Z\).

Por que escolher o declive depois de fixar a testemunha? Uma verificação
que usasse só a coordenada \(x\) confundiria \(P\) e \(-P\). Já
\(Z(P)-Z(P')=(y_P-y_{P'})-c_{\rm declive}(x_P-x_{P'})\) depende do
declive. Para pontos distintos com abscissas diferentes, há um único
declive que causa colisão; com abscissas iguais e ordenadas diferentes,
não há nenhum. Um desafio imprevisível dificulta a substituição dos lugares.

Uma avaliação isolada não é uma equivalência determinística. Uma testemunha
falsa pode coincidir num desafio particular. A solidez probabilística
depende dos limites de grau e da distribuição dos desafios, escolhidos
depois de a testemunha ficar vinculada. Esses elementos pertencem ao argumento.

\newpage
\section{Dos produtos às somas: a derivada logarítmica}

Os produtos anteriores podem ser convertidos em somas de razões. Para uma
função formal \(H(c)\ne0\), escrevemos
\[
 \frac{\mathrm d}{\mathrm dc}\log H(c)=\frac{H'(c)}{H(c)}.
\]
Esta notação não invoca logaritmos reais num corpo finito. É uma abreviatura
algébrica, cuja regra de produto segue da derivação formal:
\((UV)'/(UV)=U'/U+V'/V\).

Mantendo fixo \(c_{\rm declive}\) e derivando em
\(c=c_{\rm corte}\), os factores constantes do resultante desaparecem.
Os encontros \(B_r(c)\) movem-se formalmente com a recta. Obtemos
\begin{equation}
 \sum_{r=0}^2\frac{\mathrm d}{\mathrm dc}\log f(B_r(c))
 =\sum_j\frac{m_j}{c-Z(P_j)}.
 \label{eq:fcmp-log-derivative}
\end{equation}
Esta igualdade vale como identidade de funções racionais fora dos lugares
onde os denominadores se anulam. Mesmo que os encontros individuais
precisem de uma extensão para ser escritos como funções de \(c\), a soma
simétrica provém do resultante e está definida sobre o corpo base.

Para o divisor binário da multiplicação escalar, o lado direito é
\[
 \sum_i\frac{e_i}{c-Z(A_i)}+\frac1{c-Z(-W)},
\]
quando \(W\) é afim. Os dígitos entram linearmente; as parcelas da tabela
pública podem ser preparadas fora do circuito. A parcela de \(-W\) liga
a identidade ao ponto resultante que se pretende certificar.

A derivada elimina também a escala global: substituir \(f\) por
\(\alpha f\), com \(\alpha\ne0\), não altera \(f'/f\). Isto é
coerente com \(\operatorname{div}(\alpha f)=\operatorname{div}(f)\).
A função zero continua excluída, porque essas razões não existem para ela.

Em característica \(p\), as multiplicidades da soma são lidas no corpo,
logo módulo \(p\). Por exemplo, a derivada de uma potência \(p\)-ésima
é zero. Limites de grau muito inferiores à característica e restrições
nos dígitos impedem que este fenómeno apague a informação pretendida. A
derivada é uma ferramenta de verificação sob essas condições.

\newpage
\section{Derivar ao longo da curva}

Para calcular o lado esquerdo, não basta derivar \(f\) em \(x\) mantendo
\(y\) fixo. O ponto permanece simultaneamente na curva e na recta.
Escrevamos \(B(c)=(x(c),y(c))\), com declive fixo \(m=c_{\rm declive}\).
As equações dão
\[
 2y\frac{\mathrm dy}{\mathrm dc}
 =(3x^2+a_{\mathbb E})\frac{\mathrm dx}{\mathrm dc},
 \qquad
 \frac{\mathrm dy}{\mathrm dc}-m\frac{\mathrm dx}{\mathrm dc}=1.
\]
Resolvendo o sistema e escrevendo
\(\Delta=3x^2+a_{\mathbb E}-2my\), obtemos, quando \(\Delta\ne0\),
\[
 \frac{\mathrm dx}{\mathrm dc}=\frac{2y}{\Delta},
 \qquad
 \frac{\mathrm dy}{\mathrm dc}=\frac{3x^2+a_{\mathbb E}}{\Delta}.
\]
Se \(f_x,f_y\) são as derivadas parciais do representante reduzido,
a regra da cadeia fornece
\[
 \frac{\mathrm d}{\mathrm dc}\log f(B(c))
 =\frac{2y f_x(B)+(3x^2+a_{\mathbb E})f_y(B)}
        {f(B)(3x^2+a_{\mathbb E}-2my)}.
\]
O numerador guarda a derivação ao longo da curva; o denominador inclui
tanto o valor de \(f\) como a condição de transversalidade da recta.

Podemos conferir a identidade em pequeno. Para \(f=y-2x+9\), a recta
\(y=6\) encontra a curva em \((0,6),(7,6),(10,6)\). As três razões
calculadas pela fórmula são, respectivamente, \(14,7,2\). A soma é
\(6\) em \(\mathbb F_{17}\). Do lado dos zeros \(P,2P,-Q\), obtemos
\[
 \frac1{6-1}+\frac1{6-3}+\frac1{6-11}
 =7+6+10=6\pmod {17}.
\]
A recta \(y=1\), pelo contrário, passa pelo zero \(P\): aí uma razão
teria denominador zero. Nem toda a recta serve para este teste.

\textbf{Verificação de compreensão.} Que condição geométrica corresponde
a \(\Delta=0\), num ponto com \(y\ne0\)? O declive da recta coincide
com \((3x^2+a_{\mathbb E})/(2y)\), o declive da tangente. Os encontros
deixam de ser simples. As tangências eram úteis para construir a função;
esta fórmula de avaliação desafiada exige evitá-las.

\newpage
\section{O que as restrições realmente verificam}

Para cada encontro desafiado \(B_r\), introduz-se a razão \(o_r\) e
impõe-se uma igualdade multiplicativa:
\begin{align}
 p_{\rm den}^{(r)}o_r&=p_{\rm num}^{(r)},\nonumber\\
 p_{\rm num}^{(r)}&=2y_{B_r}f_x(B_r)
        +(3x_{B_r}^{\,2}+a_{\mathbb E})f_y(B_r),\nonumber\\
 p_{\rm den}^{(r)}&=f(B_r)
        (3x_{B_r}^{\,2}+a_{\mathbb E}-2c_{\rm declive}y_{B_r}).
 \label{eq:fcmp-divisor-cross-multiplication}
\end{align}
Esta restrição representa a divisão quando o denominador é não nulo.
Se numerador e denominador forem ambos zero, aceita qualquer \(o_r\).
A multiplicação cruzada, sozinha, não demonstra a existência da razão.

Para um denominador secreto cuja inversa deve existir, pode impor-se
explicitamente \(du=1\). Na parcela do ponto resultante,
\begin{equation}
 d_W=c_{\rm corte}+y_W+c_{\rm declive}x_W,
 \qquad d_Wd_W^{-1}=1.
 \label{eq:fcmp-output-denominator-inverse}
\end{equation}
Para denominadores públicos, o verificador pode testar a não nulidade e
calcular as inversas fora do circuito. Um protocolo que use apenas
multiplicação cruzada em certas parcelas tem de incluir os casos degenerados
na sua análise probabilística.

A identidade final compara \(\sum_r o_r\) com a soma das parcelas
prescritas pelo divisor. Porém, as avaliações de \(f,f_x,f_y\) têm de
vir dos mesmos coeficientes, já fixados, e os limites de peso têm de ser
impostos. A variável \(W\) deve satisfazer a equação da curva; os dígitos
devem satisfazer as condições da decomposição. Valores soltos que
satisfaçam uma igualdade final não bastam.

Nos argumentos do tipo Eagen, desafios imprevisíveis e limites de grau
controlam a probabilidade de uma testemunha inválida passar pelos testes.
A contagem de raízes de polinómios não nulos explica a economia: uma
diferença algébrica de grau limitado só se pode anular num conjunto
limitado de desafios. A distribuição concreta de rectas e as colisões de
\(Z\) entram igualmente nessa análise. A solidez pertence ao conjunto
das condições, não a uma única igualdade de razões.

\newpage
\section{Da função ao certificado em conhecimento zero}

Podemos finalmente reunir as peças. A afirmação é \(W=s_AA\). O provador
calcula \(W\), decompõe \(s_A\), constrói o divisor prescrito e encontra
os coeficientes de uma função não nula com o peso permitido. A função é
uma testemunha da soma nula; não precisa de ser publicada.

Para manter valores secretos, o protocolo exterior usa compromissos:
objectos públicos que vinculam o provador a valores, permitindo que esses
valores permaneçam ocultos. As suas construções e provas serão desenvolvidas
nos capítulos seguintes. Aqui importa a ordem das operações: primeiro
vincular pontos, dígitos e coeficientes; depois obter os desafios; por fim
provar as restrições das avaliações e da identidade.

O conhecimento zero não resulta apenas da presença de divisores. Publicar
os coeficientes pode revelar os zeros da função e, com eles, os pontos
seleccionados. O sistema de prova envolvente é responsável por esconder a
testemunha e por garantir que todas as restrições se referem aos mesmos
valores comprometidos.

Uma normalização exclui a função nula, por exemplo fixando a um um
coeficiente que se saiba ser não nulo. Multiplicar uma função por uma
constante não nula preserva o divisor e a derivada logarítmica. Mas escolher
de antemão um coeficiente específico pode excluir testemunhas verdadeiras
quando esse coeficiente se anula. Uma estratégia de normalização deve
tratar esse caso para assegurar a completude. O coeficiente fixado não é
necessariamente o coeficiente líder de peso elíptico do resultante.

Em FCMP++, esta maquinaria certificará multiplicações ocultas nas aberturas
de pontos e de ramos. O capítulo~\ref{chapter:fcmp-plus-plus} identificará
as relações concretas e a contagem das funções. A árvore fornece os pontos
sobre os quais se fazem alegações; as funções testemunham as multiplicações
necessárias para ligar as aberturas.

O fio matemático pode agora ser seguido sem saltos: a decomposição dá uma
soma de pontos; essa soma determina um divisor de grau zero; a função
testemunha que o divisor é principal; a reciprocidade liga a função aos
lugares prescritos; as avaliações desafiadas tornam a verificação económica.
Os limites e as condições de não nulidade conservam a validade de cada elo.
